\scanpage{299}
\mspage{40}
Alors l'application \deleted{\unsure{précéd}}
\[
 f\longmapsto\mathrm{Tr}(\deleted{\unsure{f,}}\unsure{\tau|t},f)\colon
 \operatorname{Hom}_{\underline{\mathrm{Ty}}}(\tau,\tau')\longrightarrow
 \unsure{\mathscr{P}|\mathcal{P}}_{\tau'}(\mathrm{Tr}(\deleted{\unsure{t}}\tau))
\]
est \uline{bijective} : se donner un morphisme
$f\colon\tau\unsure{\mathbin{-}|\longrightarrow}\tau'$ dans $\underline{\mathrm{Ty}}$, \uline{revient au $=$} que
de se donner une suite \inserted{str.} croissante de
flèches de transition
\[
 \deleted{\unsure{u}}\underline{u}_{i_1},\underline{u}_{i_2},\ldots\ \underline{u}_{i_{n'}}
 \qquad(n'=\operatorname{long}\tau')
\]
de $I_\tau$ $(\unsure{u_1|v_1|i_1}<i_2<\ldots<i_{n'})$, telle que
la suite des signatures \unsure{correspondante|correspondantes}
\[
 (\varepsilon_{i_1},\varepsilon_{i_2},\ldots,\varepsilon_{i_{n'}})=\tau'
\]
soit égale à $\tau'$.

\uline{Cor. 1} Pour tout $\tau$, on a des bijections
canoniques
\[
\left\{\begin{aligned}
\mathrm{Tr}^+(\tau)&\simeq\operatorname{Hom}_{\underline{\mathrm{Ty}}}(\tau,\Delta^+)\\
\mathrm{Tr}^-(\tau)&\simeq\operatorname{Hom}_{\underline{\mathrm{Ty}}}(\tau,\Delta^-)
\end{aligned}\right.,
\]
\uline{fonctorielles en $\tau$} (variant dans $\underline{\mathrm{Ty}}$).

Pour donner un sens au corollaire, il
faut définir \deleted{$\tau\mapsto\mathrm{Tr}^+(\tau)$ et $\mathrm{Tr}^-(\tau)$ comme}
\[
 \tau\unsure{\mapsto|\mathrel{\vdash}}\mathrm{Tr}^{+}(\tau),\qquad \tau\mapsto\mathrm{Tr}^{-}(\tau)
\]
comme \uline{foncteurs} contravariants en $\tau$, donc
pour un morphisme
\[
 f\colon I_{\tau_{\unsure{2|1|\,}}}\unsure{\mathbin{-}|\longrightarrow}I_{\tau_{\unsure{1|2|\,}}}
\]
dans $\mathrm{Ty}$, définir \unsure{des} applications
% archive sheet number: 298
